% =========================================================
% TIER-2 SKELETAL MANUSCRIPT INTEGRATION
% =========================================================

\subsection{Tier-2 nominal-plus-residual extension}
\label{subsec:tier2_formulation}

To move beyond the matched Tier-1 setting while preserving a disciplined controller architecture, we define a Tier-2 nominal-plus-residual extension. The Tier-2 plant is written as
\begin{equation}
\ddot s_i
=
u_i
+
\Delta_i^{\pi,\mathrm{nom}}(x)
+
r_i^\pi(x),
\label{eq:tier2_plant}
\end{equation}
where $\Delta_i^{\pi,\mathrm{nom}}(x)$ denotes the nominal interaction model retained by the controller-side evaluator, and $r_i^\pi(x)$ is a structured plant-side residual term. The controller reconstruction remains nominal:
\begin{equation}
u_i^\pi
=
v_i
-
\Delta_i^{\pi,\mathrm{nom}}(x).
\label{eq:tier2_controller}
\end{equation}

This construction preserves the controller-side flatness inversion semantics while introducing disciplined partial unmatchedness at the plant level. The central Tier-2 object is therefore the nominal-vs-plant mismatch
\begin{equation}
M_\pi(x)
:=
\Delta_{\mathrm{plant}}(x)-E_\pi^{\mathrm{nom}}(x),
\label{eq:tier2_mismatch_object}
\end{equation}
where $E_\pi^{\mathrm{nom}}(x)$ denotes the controller-side nominal evaluator and $\Delta_{\mathrm{plant}}(x)$ denotes the plant interaction actually applied in the Tier-2 dynamics. Tier-2 is thus not a new controller law; it is a controlled plant-mismatch extension of the nominal hybrid architecture.

\paragraph{Tier-2 scope.}
The Tier-2 goal is not to claim high-fidelity realism in an absolute sense. Rather, the goal is to study whether the hybrid switching story remains interpretable under structured plant-side residual mismatch while keeping the nominal controller architecture fixed.

\subsection{Tier-2 verification protocol}
\label{subsec:tier2_protocol}

The Tier-2 evidence is organized in four layers.

\paragraph{Layer 1: reduction and gap diagnostics.}
We first verify that residual-disabled Tier-2 reduces exactly to Tier-1 nominal behavior. Only after this reduction witness is established do we activate nonzero residual terms and measure nominal-vs-plant mismatch.

\paragraph{Layer 2: generic frozen-controller comparison.}
For each residual family we compare generic Tier-1 nominal and Tier-2 residual-enabled rollouts under a shared frozen-controller setup, and then repeat the comparison over paired seed sweeps.

\paragraph{Layer 3: designed scenario comparisons.}
We then construct scenario families intended to amplify mismatch-sensitive effects, including near-switch, mixed-boundary, dwell, and band-occupancy cases.

\paragraph{Layer 4: mechanism audits.}
Finally, for each residual family we audit stepwise pattern changes to determine whether mismatch acts primarily through switch-timing perturbation, candidate-history perturbation, or broader restructuring of the hybrid controller state.

\subsection{Tier-2 reduction witness}
\label{subsec:tier2_reduction}

The first Tier-2 verification step is an exact reduction witness. With the residual branch disabled, Tier-2 nominal interaction is defined to reduce exactly to the Tier-1 nominal interaction, both pointwise and at the rollout level. The reduction witness checks:
\begin{enumerate}
    \item exact pointwise equality of nominal evaluator and nominal plant interaction,
    \item zero plant-side residual,
    \item exact equality of rollout-level switching, blending, jump, and tracking series.
\end{enumerate}

The resulting witness confirms that residual-disabled Tier-2 reproduces Tier-1 nominal behavior exactly. This step is essential because it establishes that subsequent Tier-2 effects arise from the residual branch itself rather than from unintended changes to the nominal controller pathway.

\subsection{Residual family I: transverse skew}
\label{subsec:tier2_transverse_skew}

The first Tier-2 residual family is a plant-side geometry-dependent skew term, denoted \texttt{transverse\_skew}. This residual shares the Tier-1 active-edge support and introduces a small signed bias based on lateral asymmetry while keeping controller-side evaluator semantics unchanged.

\paragraph{Gap diagnostic.}
The first residual is live, structured, and nonzero. Under the nominal controller evaluator, the induced mismatch
\[
\Delta_{\mathrm{plant}}(x)-E_\pi^{\mathrm{nom}}(x)
\]
is measurable and order-sensitive through the nominal evaluator comparison, even though the plant-side interaction itself remains fixed at a given snapshot.

\paragraph{Generic rollout comparison.}
On shared frozen-controller rollouts and across paired seed sweeps, \texttt{transverse\_skew} robustly increases nominal-vs-plant mismatch burden and slightly worsens tracking and jump severity. However, it does not change
\texttt{switch\_count},
\texttt{transition\_start\_count}, or
\texttt{blend\_active\_steps}.

\paragraph{Designed mismatch-sensitive scenarios.}
Designed scenarios reveal that this first residual is not merely a generic severity perturbation. In several cases it changes switch timing, candidate histories, and effective candidate histories without changing aggregate burden counts.

\paragraph{Mechanism audit.}
The mechanism audit shows that these pattern differences are concentrated near marginal switching regions. The dominant effect is switch-timing perturbation in low-tie-margin regions, while candidate-history perturbation is present but secondary. Thus, the first residual perturbs hybrid timing structure without converting that perturbation into a change in aggregate burden under the frozen controller.

\subsection{Count-conversion attempts for the first residual}
\label{subsec:tier2_count_conversion_attempt}

Because the first residual produced timing and pattern changes without burden changes, we next constructed amplified count-conversion scenarios, including prolonged near-switch dwell, repeated grazing, and mixed-boundary dwell cases. These scenarios were intended to convert timing perturbations into changes in switch and blend counts.

The result is negative but informative: despite sustained pattern differences, the first residual still does not change aggregate switching burden. We therefore interpret the first residual as exhausted as a burden-changing mechanism under the current frozen controller.

\subsection{Residual family II: longitudinal bias}
\label{subsec:tier2_longitudinal_bias}

The second Tier-2 residual family, denoted \texttt{longitudinal\_bias}, was introduced after the first residual failed to produce burden conversion. This residual also remains plant-side only and shares Tier-1 active-edge support, but it applies a localized bias within a longitudinal band of the active interaction window.

Conceptually, this residual is intended to persist in regions where marginal switching decisions are likely to remain sensitive for multiple steps. It is therefore a stronger burden-changing candidate than a purely transverse skew.

\paragraph{Gap diagnostic.}
The second residual is live, structured, nonzero, and clearly distinct from the first residual family. It produces larger nominal-vs-plant mismatch magnitudes and substantial order-sensitive mismatch under the nominal evaluator comparison.

\paragraph{Generic rollout comparison.}
On the shared frozen rollout, \texttt{longitudinal\_bias} produces larger increases in mismatch burden, tracking degradation, and jump severity than \texttt{transverse\_skew}. The paired seed sweep confirms that these severity and mismatch increases are robust across all tested seeds. Nevertheless, aggregate switching burden remains unchanged:
\[
\Delta \texttt{switch\_count}=0,\qquad
\Delta \texttt{transition\_start\_count}=0,\qquad
\Delta \texttt{blend\_active\_steps}=0.
\]

Thus, stronger mismatch alone is still insufficient to induce burden conversion under the frozen controller.

\subsection{Designed scenarios and mechanism audit for the second residual}
\label{subsec:tier2_longitudinal_mechanism}

We next constructed designed scenario families aligned with the longitudinal-band mechanism, including longitudinal dwell, near-switch longitudinal amplification, mixed-boundary longitudinal alignment, and asymmetric band-occupancy cases.

These cases produce the strongest Tier-2 pattern consequences observed in this work. In all designed second-residual cases, we observe differences in switch timing, candidate histories, and effective candidate histories. However, these broader pattern differences still fail to change aggregate switching burden.

The mechanism audit clarifies the structure of these perturbations. Across the designed cases:
\begin{itemize}
    \item pattern-difference steps are concentrated near switch windows,
    \item tie margins at those steps remain small,
    \item mismatch norms at those steps are elevated,
    \item and pattern differences occur near the residual's longitudinal bias band.
\end{itemize}

The dominant mechanism is again marginal switch-timing perturbation. Candidate and effective-candidate perturbations are secondary, but somewhat stronger and more widespread than in the first residual family.

\subsection{Tier-2 synthesis}
\label{subsec:tier2_synthesis}

The Tier-2 results support three main conclusions.

\paragraph{First, Tier-2 is a disciplined extension of the nominal hybrid framework.}
Residual-disabled Tier-2 reduces exactly to Tier-1 nominal behavior, so Tier-2 effects are attributable to structured plant-side mismatch rather than nominal controller changes.

\paragraph{Second, structured residual mismatch has consistent closed-loop consequences.}
Across both residual families, generic rollouts and seed sweeps show robust increases in:
\begin{itemize}
    \item nominal-vs-plant mismatch burden,
    \item jump severity,
    \item and mean tracking error.
\end{itemize}

\paragraph{Third, the dominant hybrid consequence is timing/pattern sensitivity, not burden change.}
Designed scenarios and mechanism audits show that residual mismatch perturbs switch timing and candidate-pattern histories in low-margin regions. However, neither residual family converts these timing-pattern perturbations into changes in aggregate switching burden under the frozen controller.

\subsection{Tier-2 interpretation}
\label{subsec:tier2_interpretation}

Tier-2 should therefore be interpreted as a partial-unmatchedness sensitivity section, not as a new controller-improvement result. The strongest current claim is:

\begin{quote}
Under structured plant-side residual mismatch, the hybrid flatness controller remains interpretable but not insensitive: the dominant effect of mismatch is scenario-dependent perturbation of switch timing and candidate-pattern histories near low-margin switching regions, while aggregate switching burden remains unchanged under the frozen controller.
\end{quote}

This result broadens the scientific scope of the paper in an important way. Tier-1 establishes controller-relevant consequences inside a matched geometry-aware extension. Tier-2 shows that even after moving into disciplined partial unmatchedness, the hybrid structure remains analyzable: mismatch first manifests as timing and pattern sensitivity before it manifests as burden change.

\paragraph{Current limits of the Tier-2 claim.}
We do not claim that Tier-2 residuals:
\begin{itemize}
    \item change switching burden in general,
    \item justify controller redesign already,
    \item or constitute a validated external-model realism layer.
\end{itemize}
Instead, Tier-2 currently supports a narrower but stronger claim: interpretable hybrid sensitivity under structured plant-side mismatch.

\paragraph{Relationship to the Tier-1 results.}
Tier-1 remains the primary controller-consequence branch of the paper. Tier-2 complements it by showing that the hybrid mechanism persists as an interpretable object even when the plant no longer exactly matches the nominal controller evaluator.
